\chapter{Method}
\label{chap:method}

\section{The benchmark harness}
\label{sec:harness}

The harness generates synthetic documents with a controlled mix of
prose, mathematics, tables and vector figures, compiles them under each
build strategy, and records wall-clock time. \Cref{fig:pipeline} shows
the stages.

\begin{widefigure}
  \centering
  \begin{tikzpicture}[
      node distance=10mm and 11mm,
      box/.style={draw=hairline,fill=white,rounded corners=2pt,
                  minimum height=10mm,minimum width=25mm,align=center,
                  font=\sffamily\scriptsize,line width=0.6pt,inner sep=4pt},
      hot/.style={box,draw=accent,line width=1pt,text=ink},
      lbl/.style={font=\sffamily\scriptsize\itshape,text=muted,align=left},
      arr/.style={-{Stealth[length=4pt]},draw=muted,line width=0.6pt},
      feedback/.style={arr,dashed,dash pattern=on 2pt off 1.6pt}
    ]
    \node[box]               (gen)   {Document\\generator};
    \node[box,right=of gen]  (src)   {Source tree\\$n$ files};
    \node[hot,right=of src]  (build) {Build worker\\$K$ passes};
    \node[box,right=of build](pdf)   {PDF +\\timing log};

    \node[box,below=of src]   (hash)  {Content\\hashes};
    \node[box,below=of build] (cache) {Package\\cache};

    \begin{scope}[on background layer]
      \node[fill=surface,rounded corners=3pt,fit=(hash)(cache),inner sep=7pt] (panel) {};
    \end{scope}

    \draw[arr] (gen)   -- (src);
    \draw[arr] (src)   -- (build);
    \draw[arr] (build) -- (pdf);
    \draw[arr] (hash)  -- (src);
    \draw[arr] (cache) -- (build);

    % the worker writes fresh hashes back after every run
    \coordinate (turn) at ($(panel.south)+(0,-7mm)$);
    \draw[feedback] (pdf.south) |- (turn) -| (hash.south);

    \node[lbl,anchor=north] at ($(turn)+(0,-1mm)$) {new hashes written back after every run};
    \node[lbl,anchor=east,align=right] at ($(panel.west)+(-3mm,0)$) {reused\\between runs};
  \end{tikzpicture}
  \caption[Benchmark pipeline]{Stages of the benchmark harness. Only the
           build worker is timed; generation and hashing happen outside
           the measured window.}
  \label{fig:pipeline}
\end{widefigure}

\section{Build strategies}

Three strategies were compared. They differ only in what the worker is
allowed to reuse between runs.

\begin{table}[htbp]
  \centering
  \caption{The three build strategies under test.}
  \label{tab:strategies}
  \small
  \begin{tabularx}{\textwidth}{@{}l L l@{}}
    \toprule
    \textbf{\sffamily Strategy} & \textbf{\sffamily Reuses} & \textbf{\sffamily Passes $K$} \\
    \midrule
    Baseline     & nothing; a fresh container per run            & 3 \\
    Cached       & package tree, fonts, generated images          & 3 \\
    Incremental  & all of the above plus page-level content hashes & 1--3 \\
    \bottomrule
  \end{tabularx}
\end{table}

\section{Measurement protocol}

Each configuration was run \num{30} times. The first three runs are
discarded as warm-up, and we report the median of the remaining
\num{27}. Timing uses a monotonic clock taken inside the worker, so
scheduling delay in the queue is excluded.

\begin{lstlisting}[language=Python,caption={Core of the timing loop.},label={lst:timing}]
def measure(doc, strategy, repeats=30, warmup=3):
    """Return the median wall-clock build time in seconds."""
    samples = []
    for i in range(repeats):
        worker = Worker(strategy)          # fresh state for "baseline"
        t0 = time.monotonic()
        worker.build(doc)
        samples.append(time.monotonic() - t0)
    return statistics.median(samples[warmup:])
\end{lstlisting}

\begin{remark}
Using the median rather than the mean matters here: the distribution has
a long right tail caused by occasional font-cache regeneration, and a
single outlier can move the mean by more than \qty{20}{\percent}.
\end{remark}
